\section{讨论}

实验结果支持一个聚焦主张：KCH-MedRank 提供了一个可解释医学证据控制层，能够相对于强检索和语义重排序基线提升生物医学证据检索。相对于真实 MedCPT Cross-Encoder 基线，统计上最有力的比较是 Recall@10；MRR@10 和 nDCG@10 虽为正向提升，但未达到显著。因此，最稳妥的解释是 KCH-MedRank 更擅长把额外相关证据移入 top-10 证据集合，而不是保证每个问题的一位精度都更高。

消融结果表明，该方法应被理解为一个组合式生物医学重排序框架。检索和语义特征仍然是强贡献因素，而超图扩散、MeSH 层级、实体重叠和关系特征提供互补信号。这一点具有科学价值，但也约束了新颖性表述：本文贡献不是独立的超图公式，而是一种可复现地将医学知识约束注入查询级重排序的方法。

额外的选择诊断也支持保持这一克制解释。一个补充实验将完整 KCH-MedRank 的验证选择主指标改为 Recall@10，但其留出测试 Recall@10 为 0.5294，低于本文报告的以验证 MRR 为主选择的主模型 0.5332。因此，本文保留原验证选择策略，而不是在观察留出测试结果后改调目标。该诊断也说明，单纯改变选择优先级不足以稳定提分；未来提升更可能来自更好的生物医学概念归一化和更稳定的特征组选择。

PubMedQA Qwen3-8B pilot 被保留为附录式下游 RAG 系统诊断。在固定证据数量、prompt、模型和解码设置下，KCH 风格超图证据在 100 个 PubMedQA 问题上相对于 Hybrid RRF 证据提高了答案准确率和 Macro-F1。由于所有方法的引用支持率都已较高，该 pilot 不能被解读为忠实性显著改善的证明；它表明白盒证据选择可以在不训练或微调大模型的情况下影响下游 medical RAG 答案闭环，但规模和条件仍不足以支撑主要答案生成主张。

仍存在若干限制。第一，方法依赖强候选池；如果 top-100 召回较低，重排序无法恢复缺失证据。第二，PrimeKG 精确名称映射较稀疏，不能过度宣称。第三，强消融 bootstrap 检验显示完整 KCH-MedRank 与若干接近消融之间差异并不显著，因此超图、MeSH 层级、实体和关系信号应被描述为互补且可解释的信号，而不是单独占主导的因素。第四，NFCorpus 实验是 retrieval-only 且只使用检索特征，因此它支持监督重排序的稳健性，但不能证明 BioASQ 知识约束超图特征完整迁移。第五，Qwen3-8B 生成结果是 100 个 PubMedQA 问题的 pilot；在作为中心答案生成主张前，需要扩大规模并补充完整的生物医学 cross-encoder 条件。

\subsection{计算 Profile}

\begin{table}[t]
\centering
\small
\caption{增强 BioASQ 实验的计算 profile。日志记录了规模、特征维度、覆盖率和硬件信息；并未对所有组件提供可靠的端到端 wall-clock 时间。}
\label{tab:efficiency-analysis}
\begin{tabularx}{\linewidth}{p{0.22\linewidth}p{0.34\linewidth}X}
\toprule
组件 & 记录的规模或设置 & 含义 \\
\midrule
处理后数据 & 4,719 个问题；40,221 个段落；42,608 条 qrels & 适合可复现的离线重排序实验。 \\
候选池 & 每个问题 top-100；471,900 条预测记录 & 重排序受局部候选池约束，而不是受全局语料图约束。 \\
增强 KCH-MedRank 划分 & 2,832 个训练问题、944 个验证问题、943 个留出测试问题 & 超参数在留出测试之前完成选择。 \\
特征矩阵 & 45 个特征；283,200 条训练样本；377,600 条最终训练加验证样本 & 学习排序保持为查询分组的表格排序问题。 \\
局部超图范围 & 问题、top-100 段落、文档、实体、MeSH、层级和可选关系节点 & 传播成本随 top-$M$ 和局部特征覆盖率扩展。 \\
MeSH 覆盖 & 31,110 个描述符；37,356/40,221 个段落和 4,264/4,719 个问题具有 MeSH 匹配 & 受控词表特征具有较好覆盖。 \\
PrimeKG 过滤 & 扫描 8,100,498 行；保留 5,766 条项目局部关系 & 关系特征较稀疏，应保持辅助定位。 \\
稠密生物医学检索 & MedCPT dual encoder 在 RTX 4060 Laptop GPU 上使用 CUDA & GPU 加速适合稠密检索。 \\
Cross-Encoder 打分 & BioASQ CPU 完整运行在 2,400 秒后超时 & 完整 Cross-Encoder 重打分应使用 GPU 或受控子集。 \\
\bottomrule
\end{tabularx}
\end{table}


表~\ref{tab:efficiency-analysis} 总结了实验日志中可获得的计算 profile。局部超图只围绕每个问题的 top-$M$ 候选构建，而不是在全局生物医学语料上构建，因此图构建和扩散成本随候选池和局部特征覆盖率扩展。在增强 BioASQ 设置中，最终学习排序矩阵包含 45 个特征和 377,600 条训练加验证样本，使重排序保持在常规表格排序问题范围内。稠密生物医学检索受益于 GPU 加速。Cross-Encoder 重打分成本明显更高：早期通用 cross-encoder 的 BioASQ CPU 运行在 2,400 秒后超时，而最终 MedCPT Cross-Encoder 基线需要单独的受控完整测试运行。

\begin{table}[t]
\centering
\small
\caption{BioASQ 留出测试集上的重排序阶段效率对比。两种方法使用相同的增强 top-100 候选池。计时不包含第一阶段候选生成和离线 LambdaMART 训练；KCH-MedRank 包含测试时特征构建与 LightGBM 打分，MedCPT Cross-Encoder 包含分词与前向打分。}
\label{tab:reranking-efficiency}
\resizebox{\textwidth}{!}{%
\begin{tabular}{lrrrrrrl}
\toprule
方法 & 秒 & ms/问题 & 候选/秒 & Recall@10 & MRR@10 & nDCG@10 & 说明 \\
\midrule
MedCPT Cross-Encoder & 699.46 & 741.74 & 134.82 & 0.5172 & 0.7775 & 0.6390 & 分词 + 前向打分 \\
KCH-MedRank & 36.31 & 38.50 & 2597.07 & 0.5332 & 0.7882 & 0.6446 & 特征 + LightGBM \\
\bottomrule
\end{tabular}
}
\end{table}


表~\ref{tab:reranking-efficiency} 报告了 BioASQ 留出测试集上的受控重排序阶段计时。所有行使用相同的增强 top-100 候选池，共 943 个问题和 94,300 个候选段落。该比较不包含第一阶段候选生成、一次性模型加载和离线 LambdaMART 训练，并明确区分是否使用 Cross-Encoder 分数以及是否在线运行 Cross-Encoder 打分。MedCPT Cross-Encoder 在 CUDA 上完成打分需要 699.46 秒。KCH-MedRank 使用预计算的稠密生物医学语义特征，而不是 Cross-Encoder 特征；其在 CPU 上完成测试时特征构建和 LightGBM 打分需要 37.59 秒，在该计时口径下对应 18.6 倍的重排序阶段加速。no-semantic KCH 诊断在 Recall@10 上仍与 MedCPT Cross-Encoder 具有竞争力（0.5244 对 0.5172），而仅检索特征的 LambdaMART 明显更快，但 top-10 证据质量低于完整 KCH-MedRank。该结果强化了主结果的实用解释，同时避免把速度比较建立在隐藏 Cross-Encoder 特征上。包含第一阶段检索和数据加载开销的端到端生产环境基准仍留待后续工作。

未来工作应改进生物医学概念归一化，加强同义词和缩写处理，并在证据支持指标完整后扩大受控答案生成评估。更强的生成实验应衡量引用支持率、无支持声明和答案-证据实体一致性，而不仅是答案准确率。
